
\begin{obereflection}
\textbf{Latihan 4.1}
1. Gunakan kunci \texttt{Kripto} (dengan $m=7, b=10, n=26$) untuk mengenkripsi kata \texttt{halo}.
2. Diberikan cipherteks Vigen\`ere \texttt{LSimMNVWGRR} dengan kunci \texttt{soreang}, lakukan dekripsi untuk menemukan plainteksnya.
3. Buatlah matriks kunci $2 \times 2$ untuk Hill Cipher, hitung determinannya, dan tentukan apakah matriks tersebut memiliki invers modulo 26.
4. Dekripsikan cipherteks Playfair \texttt{ZB RS FY KU PG LG RK VS NL QV} dengan matriks kunci yang dibentuk dari kalimat kunci \texttt{JALAN GANESHA SEPULUH}. Tunjukkan langkah pembuangan huruf pengisi \texttt{X} pada akhir dekripsi.
5. Diketahui dua pasang plainteks--cipherteks pada affine cipher: huruf \texttt{K} menghasilkan \texttt{D}, dan huruf \texttt{T} menghasilkan \texttt{A}. Tentukan nilai $m$ dan $b$, lalu periksa apakah $\gcd(m,26)=1$.
6. Pada sebuah cipherteks Vigen\`ere ditemukan kriptogram berulang dengan jarak 15, 20, dan 30 huruf. Tentukan panjang kunci yang paling mungkin beserta alasan langkah demi langkah.
7. Enkripsikan plainteks \texttt{KRIPTOGRAFI} dengan Rail Fence cipher empat baris, lalu jelaskan mengapa cipherteks itu masih dapat dipecahkan meskipun jumlah barisnya ditambah.
\end{obereflection}
